\documentclass{article}
\usepackage{amsmath, amssymb, amsthm}
\usepackage{hyperref}
\usepackage{geometry}
\geometry{margin=1in}

\title{Erd\H{o}s Problem \#261}
\date{Last updated: 2025-08-31}
\author{Source: erdosproblems.com}

\begin{document}
\maketitle

\noindent\textbf{Status:} OPEN \quad \textbf{No prize}

\noindent\textbf{Tags:} number theory

\medskip
\noindent\textbf{Problem Statement:}

Are there infinitely many $n$ such that there exists some $t\geq 2$ and distinct integers $a_1,\ldots,a_t\geq 1$ such that\[\frac{n}{2^n}=\sum_{1\leq k\leq t}\frac{a_k}{2^{a_k}}?\]Is this true for all $n$? Is there a rational $x$ such that\[x = \sum_{k=1}^\infty \frac{a_k}{2^{a_k}}\]has at least $2^{\aleph_0}$ solutions?

\bigskip
\noindent\textbf{Remarks:}

Related to [260].

In \cite{Er88c} Erd\H{o}s notes that Cusick had a simple proof that there do exist infinitely many such $n$. Erd\H{o}s does not record what this was, but a later paper by Borwein and Loring \cite{BoLo90} provides the following proof: for every positive integer $m$ and $n=2^{m+1}-m-2$ we have\[\frac{n}{2^n}=\sum_{n<k\leq n+m}\frac{k}{2^k}.\]Tengely, Ulas, and Zygadlo \cite{TUZ20} have verified that all $n\leq 10000$ have the required property.

In \cite{Er88c} Erd\H{o}s weakens the second question to asking for the existence of a rational $x$ which has two solutions.

\bigskip
\noindent\textbf{References:}

[BoLo90] Borwein, Peter and Loring, Terry A., Some questions of {E}rd\H{o}s and {G}raham on numbers of the
form {$\sum g_n/2^{g_n}$}. Math. Comp. (1990), 377--394.
 
 [Er88c] Erd\H{o}s, P., On the irrationality of certain series: problems and results. New advances in transcendence theory (Durham, 1986) (1988), 102-109.
 
 [TUZ20] Tengely, Szabolcs and Ulas, Maciej and Zygad\l o, Jakub, On a {D}iophantine equation of {E}rd\H{o}s and {G}raham. J. Number Theory (2020), 445--459.

\bigskip
\noindent\small{Source: \url{https://www.erdosproblems.com/261}}
\end{document}
